\documentclass[aspectratio=169, 11pt]{beamer}
% ============================================================================
% Preambulo compartido para presentaciones Beamer
% Estadistica 2026-II - Pontificia Universidad Javeriana
% Incluir con: % ============================================================================
% Preambulo compartido para presentaciones Beamer
% Estadistica 2026-II - Pontificia Universidad Javeriana
% Incluir con: % ============================================================================
% Preambulo compartido para presentaciones Beamer
% Estadistica 2026-II - Pontificia Universidad Javeriana
% Incluir con: \input{../preambulo_beamer.tex} despues de \documentclass
% ============================================================================

\usepackage[T1]{fontenc}
\usepackage[utf8]{inputenc}
\usepackage[spanish]{babel}
\usepackage{amsmath, amssymb, amsfonts}
\usepackage{booktabs}
\usepackage{graphicx}
\usepackage{tikz}
\usetikzlibrary{shapes.geometric, positioning, arrows.meta, calc}
\usepackage{pgfplots}
\pgfplotsset{compat=1.18}
\usepackage{listings}
\usepackage{xcolor}
\usepackage{hyperref}
\usepackage{multicol}

% --- Tema Beamer (minimalista) ---
\usetheme{default}
\usecolortheme{default}
\usefonttheme{default}
\setbeamertemplate{navigation symbols}{}
\setbeamertemplate{caption}[numbered]
\setbeamertemplate{footline}{%
  \hspace{1em}{\tiny\url{https://github.com/polanco-jaime/Curso-Estadistica-2026-3}}\hfill\usebeamerfont{footline}\insertframenumber/\inserttotalframenumber\hspace{1em}\vspace{0.5em}%
}
\setbeamertemplate{itemize items}[circle]
\setbeamertemplate{enumerate items}[default]

% --- Colores institucionales (sutiles) ---
\definecolor{javeriana}{RGB}{0, 62, 105}
\definecolor{javerianalight}{RGB}{0, 105, 170}
\definecolor{codigobg}{RGB}{245, 245, 245}
\definecolor{codigofg}{RGB}{40, 40, 40}
\definecolor{comentario}{RGB}{100, 100, 100}
\definecolor{keyword}{RGB}{0, 100, 150}
\definecolor{string}{RGB}{180, 60, 30}

\setbeamercolor{title}{fg=javeriana}
\setbeamercolor{frametitle}{fg=javeriana}
\setbeamercolor{structure}{fg=javeriana}
\setbeamercolor{block title}{fg=white, bg=javeriana!75}
\setbeamercolor{block body}{bg=javeriana!5}
\setbeamercolor{block title alerted}{fg=white, bg=red!70!black}
\setbeamercolor{block body alerted}{bg=red!5}
\setbeamercolor{block title example}{fg=white, bg=green!40!black}
\setbeamercolor{block body example}{bg=green!5}
\setbeamercolor{item}{fg=javeriana}

\setbeamerfont{title}{series=\bfseries, size=\Large}
\setbeamerfont{frametitle}{series=\bfseries}

% --- Configuracion de listings para R ---
\lstset{
  language=R,
  basicstyle=\ttfamily\footnotesize,
  backgroundcolor=\color{codigobg},
  commentstyle=\color{comentario}\itshape,
  keywordstyle=\color{keyword}\bfseries,
  stringstyle=\color{string},
  numbers=none,
  frame=single,
  rulecolor=\color{javeriana!30},
  breaklines=true,
  showstringspaces=false,
  columns=flexible,
  morekeywords={library, ggplot, aes, geom_histogram, geom_boxplot,
    geom_point, geom_smooth, geom_density, geom_bar, geom_line,
    geom_vline, geom_hline, geom_errorbar, labs, theme_minimal,
    facet_wrap, scale_fill_manual, scale_color_manual,
    summarise, group_by, filter, mutate, select, arrange,
    read.csv, head, str, summary, mean, median, sd, var, cor,
    lm, t.test, chisq.test, wilcox.test, kruskal.test, aov,
    pnorm, qnorm, dnorm, rnorm, qt, pt, qchisq, pchisq, qf, pf,
    confint, coeftest, vcovHC, bptest, vif, cooks.distance,
    TukeyHSD, table, prop.table}
}

% --- Comandos personalizados ---
\newcommand{\R}{\texttt{R}}
\newcommand{\code}[1]{\texttt{\footnotesize #1}}
\newcommand{\variable}[1]{\texttt{#1}}
\newcommand{\paquete}[1]{\texttt{#1}}

% --- Informacion del curso ---
\institute{Pontificia Universidad Javeriana\\
  Facultad de Ciencias Econ\'omicas y Administrativas}
\date{2026-II}
 despues de \documentclass
% ============================================================================

\usepackage[T1]{fontenc}
\usepackage[utf8]{inputenc}
\usepackage[spanish]{babel}
\usepackage{amsmath, amssymb, amsfonts}
\usepackage{booktabs}
\usepackage{graphicx}
\usepackage{tikz}
\usetikzlibrary{shapes.geometric, positioning, arrows.meta, calc}
\usepackage{pgfplots}
\pgfplotsset{compat=1.18}
\usepackage{listings}
\usepackage{xcolor}
\usepackage{hyperref}
\usepackage{multicol}

% --- Tema Beamer (minimalista) ---
\usetheme{default}
\usecolortheme{default}
\usefonttheme{default}
\setbeamertemplate{navigation symbols}{}
\setbeamertemplate{caption}[numbered]
\setbeamertemplate{footline}{%
  \hspace{1em}{\tiny\url{https://github.com/polanco-jaime/Curso-Estadistica-2026-3}}\hfill\usebeamerfont{footline}\insertframenumber/\inserttotalframenumber\hspace{1em}\vspace{0.5em}%
}
\setbeamertemplate{itemize items}[circle]
\setbeamertemplate{enumerate items}[default]

% --- Colores institucionales (sutiles) ---
\definecolor{javeriana}{RGB}{0, 62, 105}
\definecolor{javerianalight}{RGB}{0, 105, 170}
\definecolor{codigobg}{RGB}{245, 245, 245}
\definecolor{codigofg}{RGB}{40, 40, 40}
\definecolor{comentario}{RGB}{100, 100, 100}
\definecolor{keyword}{RGB}{0, 100, 150}
\definecolor{string}{RGB}{180, 60, 30}

\setbeamercolor{title}{fg=javeriana}
\setbeamercolor{frametitle}{fg=javeriana}
\setbeamercolor{structure}{fg=javeriana}
\setbeamercolor{block title}{fg=white, bg=javeriana!75}
\setbeamercolor{block body}{bg=javeriana!5}
\setbeamercolor{block title alerted}{fg=white, bg=red!70!black}
\setbeamercolor{block body alerted}{bg=red!5}
\setbeamercolor{block title example}{fg=white, bg=green!40!black}
\setbeamercolor{block body example}{bg=green!5}
\setbeamercolor{item}{fg=javeriana}

\setbeamerfont{title}{series=\bfseries, size=\Large}
\setbeamerfont{frametitle}{series=\bfseries}

% --- Configuracion de listings para R ---
\lstset{
  language=R,
  basicstyle=\ttfamily\footnotesize,
  backgroundcolor=\color{codigobg},
  commentstyle=\color{comentario}\itshape,
  keywordstyle=\color{keyword}\bfseries,
  stringstyle=\color{string},
  numbers=none,
  frame=single,
  rulecolor=\color{javeriana!30},
  breaklines=true,
  showstringspaces=false,
  columns=flexible,
  morekeywords={library, ggplot, aes, geom_histogram, geom_boxplot,
    geom_point, geom_smooth, geom_density, geom_bar, geom_line,
    geom_vline, geom_hline, geom_errorbar, labs, theme_minimal,
    facet_wrap, scale_fill_manual, scale_color_manual,
    summarise, group_by, filter, mutate, select, arrange,
    read.csv, head, str, summary, mean, median, sd, var, cor,
    lm, t.test, chisq.test, wilcox.test, kruskal.test, aov,
    pnorm, qnorm, dnorm, rnorm, qt, pt, qchisq, pchisq, qf, pf,
    confint, coeftest, vcovHC, bptest, vif, cooks.distance,
    TukeyHSD, table, prop.table}
}

% --- Comandos personalizados ---
\newcommand{\R}{\texttt{R}}
\newcommand{\code}[1]{\texttt{\footnotesize #1}}
\newcommand{\variable}[1]{\texttt{#1}}
\newcommand{\paquete}[1]{\texttt{#1}}

% --- Informacion del curso ---
\institute{Pontificia Universidad Javeriana\\
  Facultad de Ciencias Econ\'omicas y Administrativas}
\date{2026-II}
 despues de \documentclass
% ============================================================================

\usepackage[T1]{fontenc}
\usepackage[utf8]{inputenc}
\usepackage[spanish]{babel}
\usepackage{amsmath, amssymb, amsfonts}
\usepackage{booktabs}
\usepackage{graphicx}
\usepackage{tikz}
\usetikzlibrary{shapes.geometric, positioning, arrows.meta, calc}
\usepackage{pgfplots}
\pgfplotsset{compat=1.18}
\usepackage{listings}
\usepackage{xcolor}
\usepackage{hyperref}
\usepackage{multicol}

% --- Tema Beamer (minimalista) ---
\usetheme{default}
\usecolortheme{default}
\usefonttheme{default}
\setbeamertemplate{navigation symbols}{}
\setbeamertemplate{caption}[numbered]
\setbeamertemplate{footline}{%
  \hspace{1em}{\tiny\url{https://github.com/polanco-jaime/Curso-Estadistica-2026-3}}\hfill\usebeamerfont{footline}\insertframenumber/\inserttotalframenumber\hspace{1em}\vspace{0.5em}%
}
\setbeamertemplate{itemize items}[circle]
\setbeamertemplate{enumerate items}[default]

% --- Colores institucionales (sutiles) ---
\definecolor{javeriana}{RGB}{0, 62, 105}
\definecolor{javerianalight}{RGB}{0, 105, 170}
\definecolor{codigobg}{RGB}{245, 245, 245}
\definecolor{codigofg}{RGB}{40, 40, 40}
\definecolor{comentario}{RGB}{100, 100, 100}
\definecolor{keyword}{RGB}{0, 100, 150}
\definecolor{string}{RGB}{180, 60, 30}

\setbeamercolor{title}{fg=javeriana}
\setbeamercolor{frametitle}{fg=javeriana}
\setbeamercolor{structure}{fg=javeriana}
\setbeamercolor{block title}{fg=white, bg=javeriana!75}
\setbeamercolor{block body}{bg=javeriana!5}
\setbeamercolor{block title alerted}{fg=white, bg=red!70!black}
\setbeamercolor{block body alerted}{bg=red!5}
\setbeamercolor{block title example}{fg=white, bg=green!40!black}
\setbeamercolor{block body example}{bg=green!5}
\setbeamercolor{item}{fg=javeriana}

\setbeamerfont{title}{series=\bfseries, size=\Large}
\setbeamerfont{frametitle}{series=\bfseries}

% --- Configuracion de listings para R ---
\lstset{
  language=R,
  basicstyle=\ttfamily\footnotesize,
  backgroundcolor=\color{codigobg},
  commentstyle=\color{comentario}\itshape,
  keywordstyle=\color{keyword}\bfseries,
  stringstyle=\color{string},
  numbers=none,
  frame=single,
  rulecolor=\color{javeriana!30},
  breaklines=true,
  showstringspaces=false,
  columns=flexible,
  morekeywords={library, ggplot, aes, geom_histogram, geom_boxplot,
    geom_point, geom_smooth, geom_density, geom_bar, geom_line,
    geom_vline, geom_hline, geom_errorbar, labs, theme_minimal,
    facet_wrap, scale_fill_manual, scale_color_manual,
    summarise, group_by, filter, mutate, select, arrange,
    read.csv, head, str, summary, mean, median, sd, var, cor,
    lm, t.test, chisq.test, wilcox.test, kruskal.test, aov,
    pnorm, qnorm, dnorm, rnorm, qt, pt, qchisq, pchisq, qf, pf,
    confint, coeftest, vcovHC, bptest, vif, cooks.distance,
    TukeyHSD, table, prop.table}
}

% --- Comandos personalizados ---
\newcommand{\R}{\texttt{R}}
\newcommand{\code}[1]{\texttt{\footnotesize #1}}
\newcommand{\variable}[1]{\texttt{#1}}
\newcommand{\paquete}[1]{\texttt{#1}}

% --- Informacion del curso ---
\institute{Pontificia Universidad Javeriana\\
  Facultad de Ciencias Econ\'omicas y Administrativas}
\date{2026-II}

\title{Sesión 2: Estadística Descriptiva Gráfica, Correlación y EDA}
\subtitle{Visualización de datos, correlación y Paradoja de Simpson}
\author{Prof.\ Jaime Polanco Jim\'enez}
\date{Vie 14 ago 2026}
\begin{document}

\begin{frame}
\titlepage
\end{frame}

\begin{frame}{Agenda}
\tableofcontents
\end{frame}

% ============================================================================
\section{Distribuciones de frecuencia}
% ============================================================================

\begin{frame}{Distribución de frecuencias}
\begin{block}{Definición}
Una \textbf{distribución de frecuencias} es una tabla que organiza los datos en clases o categor\'ias, mostrando cu\'antas observaciones caen en cada una.
\end{block}

\textbf{Ejemplo}: si encuestan a 30 compa\~neros sobre cu\'antas horas de TikTok ven al d\'ia, la tabla de frecuencias organiza esas respuestas en categor\'ias (0--1h, 1--2h, 2--3h, ...) y cuenta cu\'antos caen en cada una.

\textbf{Tipos de frecuencia}:
\begin{itemize}
\item \textbf{Frecuencia absoluta} ($f_i$): número de observaciones en la clase $i$
\item \textbf{Frecuencia relativa} ($h_i$): proporción, $h_i = f_i / n$
\item \textbf{Frecuencia acumulada} ($F_i$): suma de frecuencias hasta la clase $i$
\item \textbf{Frecuencia relativa acumulada} ($H_i$): proporción acumulada hasta $i$
\end{itemize}
\end{frame}

\begin{frame}{Ejemplo: tabla de frecuencias}
\begin{exampleblock}{Puntajes de inglés agrupados}
\centering
\small
\begin{tabular}{lcccc}
\toprule
Clase & $f_i$ & $h_i$ & $F_i$ & $H_i$ \\
\midrule
100--119 & 45 & 0.045 & 45 & 0.045 \\
120--139 & 180 & 0.180 & 225 & 0.225 \\
140--159 & 420 & 0.420 & 645 & 0.645 \\
160--179 & 280 & 0.280 & 925 & 0.925 \\
180--200 & 75 & 0.075 & 1000 & 1.000 \\
\bottomrule
\end{tabular}
\end{exampleblock}

\textbf{Lectura}: el 42\% de los estudiantes obtiene entre 140 y 159 puntos; el 64.5\% est\'a por debajo de 160. La misma l\'ogica aplica si agrupan el screen time diario del sal\'on en intervalos.
\end{frame}

\begin{frame}[fragile]{Tablas de frecuencia en R}
\begin{lstlisting}
# Para variables categoricas: usar table()
tipo_ies <- c("Publica", "Privada", "Privada", "Publica",
              "Privada", "Privada", "Publica", "Privada")

# Frecuencias absolutas
table(tipo_ies)
# Privada  Publica
#       5        3

# Frecuencias relativas
prop.table(table(tipo_ies))
# Privada  Publica
#   0.625    0.375

# Con datos ICFES
freq_ies <- table(icfes$INST_CARACTER_ACADEMICO)
prop.table(freq_ies)
\end{lstlisting}
\end{frame}

\begin{frame}[fragile]{Distribución de frecuencias agrupadas (1/2)}
\begin{lstlisting}
# Para variables continuas: usar cut() para crear intervalos
puntajes <- icfes$MOD_INGLES_PUNT

# Crear intervalos
intervalos <- cut(puntajes,
                  breaks = seq(100, 200, by = 20),
                  include.lowest = TRUE,
                  right = FALSE)

# Tabla de frecuencias
freq_tabla <- table(intervalos)
freq_tabla

# Frecuencias relativas
prop.table(freq_tabla)
\end{lstlisting}
\end{frame}

\begin{frame}[fragile]{Distribución de frecuencias agrupadas (2/2)}
\begin{lstlisting}
# Tabla completa con frecuencias acumuladas
library(tidyverse)
tibble(
  clase = names(freq_tabla),
  f = as.numeric(freq_tabla),
  h = f / sum(f),
  F = cumsum(f),
  H = cumsum(h)
)
\end{lstlisting}
\end{frame}

% ============================================================================
\section{Visualización de una variable}
% ============================================================================

\begin{frame}{Histograma}
\begin{block}{Definición}
Un \textbf{histograma} muestra la distribuci\'on de frecuencias de una variable continua mediante barras contiguas. Imaginen un histograma del screen time diario de todos en el sal\'on: cada barra muestra cu\'antos caen en un rango de horas.
\end{block}

\begin{columns}[T]
\begin{column}{0.48\textwidth}
\textbf{Componentes}:
\begin{itemize}
\item \textbf{Bins} (intervalos): rangos en el eje x
\item \textbf{Altura}: frecuencia o densidad
\item Área total = 1 (con densidad)
\end{itemize}
\end{column}
\begin{column}{0.48\textwidth}
\textbf{Interpretación}:
\begin{itemize}
\item Centro: ¿dónde se concentran?
\item Dispersión: ¿qué tan amplia?
\item Forma: ¿simétrica, sesgada?
\item Outliers: ¿valores aislados?
\end{itemize}
\end{column}
\end{columns}

\vspace{0.3em}
\begin{alertblock}{Número de bins}
Demasiados $\to$ ruido; muy pocos $\to$ se pierde detalle.
\textbf{Regla de Sturges}: $k = 1 + 3.322 \log_{10}(n)$
\end{alertblock}
\end{frame}

\begin{frame}[fragile]{Histograma en R con ggplot2}
\begin{lstlisting}
library(ggplot2)

# Histograma basico
ggplot(icfes, aes(x = MOD_INGLES_PUNT)) +
  geom_histogram(bins = 30, fill = "steelblue", color = "white") +
  labs(title = "Distribucion de puntajes de ingles - Saber Pro NI",
       x = "Puntaje de ingles",
       y = "Frecuencia") +
  theme_minimal()

# Histograma con densidad
ggplot(icfes, aes(x = MOD_INGLES_PUNT)) +
  geom_histogram(aes(y = after_stat(density)), bins = 30,
                 fill = "steelblue", color = "white", alpha = 0.7) +
  geom_density(color = "red", linewidth = 1) +
  labs(title = "Distribucion de puntajes de ingles",
       x = "Puntaje", y = "Densidad") +
  theme_minimal()
\end{lstlisting}
\end{frame}

\begin{frame}[fragile]{Estimación de densidad kernel (KDE)}
\begin{block}{Kernel Density Estimation}
Alternativa suavizada al histograma. Estima la función de densidad de probabilidad subyacente.
\end{block}

\begin{lstlisting}
# Density plot
ggplot(icfes, aes(x = MOD_INGLES_PUNT)) +
  geom_density(fill = "lightblue", alpha = 0.5, color = "darkblue") +
  geom_vline(aes(xintercept = mean(MOD_INGLES_PUNT, na.rm = TRUE)),
             color = "red", linetype = "dashed", linewidth = 1) +
  geom_vline(aes(xintercept = median(MOD_INGLES_PUNT, na.rm = TRUE)),
             color = "green", linetype = "dashed", linewidth = 1) +
  labs(title = "Estimacion de densidad - puntaje de ingles",
       x = "Puntaje", y = "Densidad",
       caption = "Rojo = media, Verde = mediana") +
  theme_minimal()
\end{lstlisting}

\textbf{Ventajas}: suave, fácil de interpretar\\
\textbf{Desventajas}: sensible al parámetro de ancho de banda
\end{frame}

\begin{frame}{Box plot (diagrama de caja)}
\begin{block}{Definición}
Un \textbf{box plot} es un gráfico que muestra el five-number summary y los outliers.
\end{block}

\textbf{Componentes}:
\begin{itemize}
\item \textbf{Caja}: va desde Q1 hasta Q3 (contiene el 50\% central)
\item \textbf{Línea en la caja}: mediana (Q2)
\item \textbf{Bigotes}: se extienden hasta el valor más lejano dentro de 1.5 $\times$ IQR
\item \textbf{Puntos}: outliers (fuera de los bigotes)
\end{itemize}

\begin{center}
\begin{tikzpicture}[scale=0.8]
\draw[thick] (0,0) rectangle (4,1.5);
\draw[thick, red] (2.5,0) -- (2.5,1.5);
\draw[thick] (-1.5,0.75) -- (0,0.75);
\draw[thick] (4,0.75) -- (5.5,0.75);
\draw[thick] (-1.5,0.5) -- (-1.5,1);
\draw[thick] (5.5,0.5) -- (5.5,1);
\fill (6.5,0.75) circle (0.1);
\fill (7,0.75) circle (0.1);
\node[below] at (-1.5,-0.3) {Min};
\node[below] at (0,-0.3) {Q1};
\node[below] at (2.5,-0.3) {Q2};
\node[below] at (4,-0.3) {Q3};
\node[below] at (5.5,-0.3) {Max};
\node[below] at (6.75,-0.3) {Outliers};
\draw[<->] (0,-0.8) -- (4,-0.8);
\node[below] at (2,-0.8) {IQR};
\end{tikzpicture}
\end{center}

\textbf{Ventajas}: compacto, f\'acil comparar grupos, identifica outliers.
Ej: comparar horas de sue\~no entre los que usan TikTok antes de dormir vs.\ los que no.
\end{frame}

\begin{frame}[fragile]{Box plot en R (1/2)}
\begin{lstlisting}
# Box plot basico (vertical)
ggplot(icfes, aes(y = MOD_INGLES_PUNT)) +
  geom_boxplot(fill = "lightblue", outlier.color = "red") +
  labs(title = "Box plot de puntajes de ingles",
       y = "Puntaje de ingles") +
  theme_minimal()

# Box plot horizontal
ggplot(icfes, aes(x = MOD_INGLES_PUNT)) +
  geom_boxplot(fill = "lightgreen", outlier.color = "red") +
  labs(title = "Puntajes de ingles - NI 2024",
       x = "Puntaje de ingles") +
  theme_minimal()
\end{lstlisting}
\end{frame}

\begin{frame}[fragile]{Box plot en R (2/2): identificar outliers}
\begin{lstlisting}
# Identificar outliers numericamente con metodo IQR
Q1 <- quantile(icfes$MOD_INGLES_PUNT, 0.25, na.rm = TRUE)
Q3 <- quantile(icfes$MOD_INGLES_PUNT, 0.75, na.rm = TRUE)
IQR_val <- Q3 - Q1

outliers <- icfes %>%
  filter(MOD_INGLES_PUNT < Q1 - 1.5*IQR_val |
         MOD_INGLES_PUNT > Q3 + 1.5*IQR_val)

cat("Numero de outliers:", nrow(outliers))
\end{lstlisting}
\end{frame}

\begin{frame}[fragile]{Box plot comparativo: por tipo de IES}
\begin{lstlisting}
# Comparar por tipo de IES
ggplot(icfes, aes(x = INST_CARACTER_ACADEMICO,
                  y = MOD_INGLES_PUNT,
                  fill = INST_CARACTER_ACADEMICO)) +
  geom_boxplot(outlier.color = "red", outlier.alpha = 0.5) +
  labs(title = "Puntajes de ingles por tipo de IES",
       x = "Tipo de institucion",
       y = "Puntaje de ingles") +
  theme_minimal() +
  theme(legend.position = "none",
        axis.text.x = element_text(angle = 45, hjust = 1))
\end{lstlisting}
\end{frame}

\begin{frame}[fragile]{Box plot comparativo: por estrato}
\begin{lstlisting}
# Comparar puntaje de ingles por estrato socioeconomico
icfes %>%
  filter(!is.na(FAMI_ESTRATOVIVIENDA)) %>%
  ggplot(aes(x = factor(FAMI_ESTRATOVIVIENDA),
             y = MOD_INGLES_PUNT,
             fill = factor(FAMI_ESTRATOVIVIENDA))) +
  geom_boxplot() +
  labs(title = "Puntajes de ingles por estrato",
       x = "Estrato", y = "Puntaje de ingles") +
  theme_minimal() +
  theme(legend.position = "none")
\end{lstlisting}
\end{frame}

\begin{frame}[fragile]{Gráfico de barras: frecuencias}
\begin{lstlisting}
# Grafico de barras simple
ggplot(icfes, aes(x = INST_CARACTER_ACADEMICO)) +
  geom_bar(fill = "steelblue") +
  labs(title = "Estudiantes NI por tipo de IES",
       x = "Tipo de institucion", y = "Frecuencia") +
  theme_minimal() +
  theme(axis.text.x = element_text(angle = 45, hjust = 1))
\end{lstlisting}
\end{frame}

\begin{frame}[fragile]{Gráfico de barras: proporciones}
\begin{lstlisting}
# Con proporciones y etiquetas
icfes %>%
  count(INST_CARACTER_ACADEMICO) %>%
  mutate(prop = n / sum(n)) %>%
  ggplot(aes(x = reorder(INST_CARACTER_ACADEMICO, -prop),
             y = prop)) +
  geom_bar(stat = "identity", fill = "coral") +
  geom_text(aes(label = paste0(round(prop*100, 1), "%")),
            vjust = -0.5) +
  labs(title = "Proporcion de estudiantes por tipo de IES",
       x = "Tipo de institucion", y = "Proporcion") +
  theme_minimal() +
  theme(axis.text.x = element_text(angle = 45, hjust = 1))
\end{lstlisting}
\end{frame}

\begin{frame}{Buenas prácticas de visualización}
\begin{alertblock}{Principios de visualización efectiva}
\begin{enumerate}
\item \textbf{Claridad}: título descriptivo, ejes etiquetados con unidades
\item \textbf{Simplicidad}: eliminar elementos innecesarios (chartjunk)
\item \textbf{Honestidad}: ejes que empiezan en cero (para barras), escalas apropiadas
\item \textbf{Accesibilidad}: paletas de colores amigables para daltónicos
\item \textbf{Contexto}: incluir leyenda, fuente de datos, notas explicativas
\end{enumerate}
\end{alertblock}

\vspace{0.5em}
\textbf{Errores comunes a evitar}:
\begin{itemize}
\item Gráficos 3D innecesarios (distorsionan percepción)
\item Demasiados colores (usar 3-5 máximo)
\item Ejes truncados o escalas manipuladas para exagerar diferencias (esos gr\'aficos enga\~nosos que ven en Twitter/X con ejes manipulados)
\item Gráficos de pastel con muchas categorías (preferir barras)
\item Falta de contexto o unidades
\end{itemize}
\end{frame}

% ============================================================================
\section{Visualización de dos variables}
% ============================================================================

\begin{frame}{Scatter plot (diagrama de dispersión)}
\begin{block}{Definición}
Un \textbf{scatter plot} muestra la relaci\'on entre dos variables cuantitativas, con cada punto representando una observaci\'on. Ej: horas de estudio vs.\ nota del parcial --- cada punto es un compa\~nero del sal\'on.
\end{block}

\textbf{Qu\'e buscar}:
\begin{itemize}
\item \textbf{Dirección}: positiva, negativa o ninguna
\item \textbf{Forma}: lineal, curva, sin patrón
\item \textbf{Fuerza}: puntos cerca de una línea (fuerte) o dispersos (débil)
\item \textbf{Outliers}: puntos alejados del patrón general
\end{itemize}
\end{frame}

\begin{frame}{Scatter plot: ejemplo de interpretación}
\begin{exampleblock}{Saber 11 inglés (eje x) vs Saber Pro inglés (eje y)}
\begin{itemize}
\item Dirección positiva: mejor Saber 11 $\to$ mejor Saber Pro
\item Forma aproximadamente lineal
\item Fuerza moderada a fuerte ($r \approx 0.6$--$0.7$)
\end{itemize}
\end{exampleblock}

\vspace{0.5em}
La correlación cuantifica esta relación (ver sección siguiente).
\end{frame}

\begin{frame}[fragile]{Scatter plot en R}
\begin{lstlisting}
# Scatter plot basico
ggplot(icfes, aes(x = PUNT_INGLES_S11, y = MOD_INGLES_PUNT)) +
  geom_point(alpha = 0.3, color = "steelblue") +
  labs(title = "Relacion entre Saber 11 y Saber Pro (Ingles)",
       x = "Puntaje Saber 11 Ingles",
       y = "Puntaje Saber Pro Ingles") +
  theme_minimal()

# Con linea de tendencia
ggplot(icfes, aes(x = PUNT_INGLES_S11, y = MOD_INGLES_PUNT)) +
  geom_point(alpha = 0.3, color = "steelblue") +
  geom_smooth(method = "lm", color = "red", se = TRUE) +
  labs(title = "Relacion Saber 11 vs Saber Pro (Ingles)",
       x = "Saber 11 Ingles",
       y = "Saber Pro Ingles",
       caption = paste("r =",
         round(cor(icfes$PUNT_INGLES_S11, icfes$MOD_INGLES_PUNT,
                   use = "complete.obs"), 3))) +
  theme_minimal()
\end{lstlisting}
\end{frame}

\begin{frame}[fragile]{Scatter plot: colorear por grupo}
\begin{lstlisting}
# Colorear por tipo de IES
ggplot(icfes, aes(x = PUNT_INGLES_S11,
                  y = MOD_INGLES_PUNT,
                  color = INST_CARACTER_ACADEMICO)) +
  geom_point(alpha = 0.4) +
  geom_smooth(method = "lm", se = FALSE) +
  labs(title = "S11 vs SPro por tipo de IES",
       x = "Saber 11 Ingles",
       y = "Saber Pro Ingles",
       color = "Tipo de IES") +
  theme_minimal()
\end{lstlisting}
\end{frame}

\begin{frame}[fragile]{Scatter plot: facetas por estrato}
\begin{lstlisting}
# Facetas por estrato socioeconomico
icfes %>%
  filter(!is.na(FAMI_ESTRATOVIVIENDA),
         FAMI_ESTRATOVIVIENDA %in% c(1,2,3,4,5,6)) %>%
  ggplot(aes(x = PUNT_INGLES_S11, y = MOD_INGLES_PUNT)) +
  geom_point(alpha = 0.3, color = "steelblue") +
  geom_smooth(method = "lm", color = "red", se = FALSE) +
  facet_wrap(~ FAMI_ESTRATOVIVIENDA, nrow = 2) +
  labs(title = "Relacion S11 vs SPro por estrato",
       x = "Saber 11 Ingles", y = "Saber Pro Ingles") +
  theme_minimal()
\end{lstlisting}
\end{frame}

\begin{frame}[fragile]{Heatmap: puntaje promedio por departamento (1/2)}
\begin{lstlisting}
# Calcular puntaje promedio por departamento
mapa_depto <- icfes %>%
  filter(!is.na(ESTU_DEPTO_RESIDE)) %>%
  group_by(ESTU_DEPTO_RESIDE) %>%
  summarise(
    n = n(),
    puntaje_medio = mean(MOD_INGLES_PUNT, na.rm = TRUE)
  ) %>%
  arrange(desc(puntaje_medio))
\end{lstlisting}
\end{frame}

\begin{frame}[fragile]{Heatmap: puntaje promedio por departamento (2/2)}
\begin{lstlisting}
# Heatmap simple (solo tabla)
ggplot(mapa_depto %>% top_n(20, n),
       aes(x = reorder(ESTU_DEPTO_RESIDE, puntaje_medio),
           y = 1, fill = puntaje_medio)) +
  geom_tile(color = "white") +
  geom_text(aes(label = round(puntaje_medio, 1)), color = "white") +
  scale_fill_gradient(low = "red", high = "green") +
  labs(title = "Puntaje promedio de ingles por departamento (Top 20)",
       x = "Departamento", fill = "Puntaje medio") +
  theme_minimal() +
  theme(axis.text.x = element_text(angle = 45, hjust = 1),
        axis.title.y = element_blank(),
        axis.text.y = element_blank())
\end{lstlisting}
\end{frame}

\begin{frame}[fragile]{Barras lado a lado (dodge)}
\begin{lstlisting}
# Preparar datos: contar por tipo de IES y estrato
datos_barras <- icfes %>%
  filter(!is.na(FAMI_ESTRATOVIVIENDA),
         FAMI_ESTRATOVIVIENDA %in% c(2,3,4,5)) %>%
  count(INST_CARACTER_ACADEMICO, FAMI_ESTRATOVIVIENDA)

# Barras lado a lado
ggplot(datos_barras, aes(x = factor(FAMI_ESTRATOVIVIENDA),
                         y = n,
                         fill = INST_CARACTER_ACADEMICO)) +
  geom_bar(stat = "identity", position = "dodge") +
  labs(title = "Estudiantes por estrato y tipo IES",
       x = "Estrato", y = "Frecuencia", fill = "Tipo IES") +
  theme_minimal()
\end{lstlisting}
\end{frame}

\begin{frame}[fragile]{Barras apiladas (100\%)}
\begin{lstlisting}
# Barras apiladas proporcionales
ggplot(datos_barras, aes(x = factor(FAMI_ESTRATOVIVIENDA),
                         y = n,
                         fill = INST_CARACTER_ACADEMICO)) +
  geom_bar(stat = "identity", position = "fill") +
  labs(title = "Composicion por tipo IES segun estrato",
       x = "Estrato", y = "Proporcion", fill = "Tipo IES") +
  theme_minimal()
\end{lstlisting}

\texttt{position = "dodge"} compara magnitudes; \texttt{position = "fill"} compara composición relativa.
\end{frame}

\begin{frame}[fragile]{Tablas cruzadas con \code{table()}}
\begin{lstlisting}
# Tabla cruzada de dos variables categoricas
tabla_cruzada <- table(icfes$INST_CARACTER_ACADEMICO,
                       icfes$FAMI_ESTRATOVIVIENDA)
tabla_cruzada

# Proporciones por fila (suman 1 en cada fila)
prop.table(tabla_cruzada, margin = 1)

# Proporciones por columna (suman 1 en cada columna)
prop.table(tabla_cruzada, margin = 2)
\end{lstlisting}
\end{frame}

\begin{frame}[fragile]{Tablas cruzadas con \paquete{tidyverse}}
\begin{lstlisting}
# Proporciones por tipo de IES con tidyverse
icfes %>%
  filter(!is.na(FAMI_ESTRATOVIVIENDA)) %>%
  count(INST_CARACTER_ACADEMICO, FAMI_ESTRATOVIVIENDA) %>%
  group_by(INST_CARACTER_ACADEMICO) %>%
  mutate(proporcion = n / sum(n))

# Prueba chi-cuadrado de asociacion
chisq.test(tabla_cruzada)
\end{lstlisting}

Veremos pruebas de hipótesis formales en sesiones futuras.
\end{frame}

% ============================================================================
\section{Covarianza y correlación}
% ============================================================================

\begin{frame}{Covarianza}
\begin{block}{Definición}
La \textbf{covarianza} mide la variación conjunta de dos variables. Indica si tienden a variar juntas en la misma dirección o en direcciones opuestas.
\end{block}

\textbf{Covarianza muestral}:
\[
s_{xy} = \frac{1}{n-1} \sum_{i=1}^n (x_i - \bar{x})(y_i - \bar{y})
\]

\textbf{Interpretación}:
\begin{itemize}
\item $s_{xy} > 0$: relación positiva (cuando x aumenta, y tiende a aumentar)
\item $s_{xy} < 0$: relación negativa (cuando x aumenta, y tiende a disminuir)
\item $s_{xy} \approx 0$: no hay relaci\'on lineal
\end{itemize}

\textbf{Ejemplo}: los likes y comentarios en un post de Instagram tienden a subir juntos $\to$ covarianza positiva.

\begin{alertblock}{Limitaci\'on}
La magnitud de $s_{xy}$ depende de las unidades de x e y. No se puede comparar entre pares de variables diferentes. \alert{Solución: usar correlación}.
\end{alertblock}
\end{frame}

\begin{frame}{Correlación de Pearson}
\begin{block}{Definición}
El \textbf{coeficiente de correlación de Pearson} (r) es la covarianza estandarizada. Mide la fuerza y dirección de la relación \alert{lineal} entre dos variables.
\[
r = \frac{s_{xy}}{s_x \cdot s_y} = \frac{\sum(x_i - \bar{x})(y_i - \bar{y})}{\sqrt{\sum(x_i-\bar{x})^2 \sum(y_i-\bar{y})^2}}
\]
\end{block}

\textbf{Propiedades}:
\begin{itemize}
\item Rango: $-1 \leq r \leq 1$
\item $r = 1$: relación lineal positiva perfecta
\item $r = -1$: relación lineal negativa perfecta
\item $r = 0$: no hay relación lineal
\item Sin unidades (adimensional)
\item Simétrico: $r_{xy} = r_{yx}$
\end{itemize}
\end{frame}

\begin{frame}{Interpretación de la correlación}
\begin{table}
\centering
\begin{tabular}{ll}
\toprule
\textbf{Valor de $|r|$} & \textbf{Interpretación} \\
\midrule
$0.0 - 0.3$ & Débil o nula \\
$0.3 - 0.5$ & Moderada débil \\
$0.5 - 0.7$ & Moderada \\
$0.7 - 0.9$ & Fuerte \\
$0.9 - 1.0$ & Muy fuerte \\
\bottomrule
\end{tabular}
\end{table}

\vspace{0.3em}
\textbf{Nota}: estos rangos son gu\'ias orientativas, no reglas absolutas.

\textbf{Ejemplo}: $r = 0.8$ entre horas de estudio y nota del parcial indica una relaci\'on fuerte positiva: a m\'as horas, mejor nota (en promedio).
\end{frame}

\begin{frame}{Advertencias sobre la correlación}
\begin{alertblock}{Tres advertencias clave}
\begin{enumerate}
\item \alert{Solo mide relaciones LINEALES}. Puede haber relaciones no lineales fuertes con $r \approx 0$.
\item \alert{Correlación $\neq$ causalidad}. Puede deberse a variables confusoras, causalidad inversa o coincidencia.
\item \alert{Outliers} pueden inflar o deflacionar $r$ artificialmente.
\end{enumerate}
\end{alertblock}

\vspace{0.5em}
\textbf{Ejemplo}: la correlaci\'on entre consumo de caf\'e y buenas notas es alta, pero la variable confusora es que los que toman m\'as caf\'e tambi\'en estudian m\'as horas.
\end{frame}

\begin{frame}[fragile]{Correlación en R (1/2)}
\begin{lstlisting}
# Correlacion entre dos variables
cor(icfes$PUNT_INGLES_S11, icfes$MOD_INGLES_PUNT,
    use = "complete.obs")
# [1] 0.682
# use = "complete.obs" elimina pares con NA

# Matriz de correlacion de multiples variables
variables_interes <- icfes %>%
  select(MOD_INGLES_PUNT, MOD_LECTURA_CRITICA_PUNT,
         MOD_RAZONA_CUANTITAT_PUNT, MOD_COMPETEN_CIUDADA_PUNT,
         MOD_COMUNICACION_ESCRITA_PUNT)

matriz_cor <- cor(variables_interes, use = "complete.obs")
round(matriz_cor, 3)
\end{lstlisting}
\end{frame}

\begin{frame}[fragile]{Correlación en R (2/2): visualización}
\begin{lstlisting}
# Visualizar matriz de correlacion con corrplot
library(corrplot)
corrplot(matriz_cor, method = "color", type = "upper",
         addCoef.col = "black", tl.col = "black",
         tl.srt = 45, diag = FALSE)
\end{lstlisting}

\texttt{corrplot} permite ver patrones rápidamente: azul = positiva, rojo = negativa, intensidad = fuerza.
\end{frame}

\begin{frame}{Correlación de Spearman}
\begin{block}{Definición}
La \textbf{correlación de Spearman} ($\rho$, rho) es la correlación de Pearson aplicada a los \alert{rangos} de los datos, no a los valores originales.
\end{block}

\textbf{Cuándo usar Spearman}:
\begin{itemize}
\item Relación monotónica pero no necesariamente lineal
\item Datos ordinales
\item Presencia de outliers (es más robusta)
\item Distribuciones muy sesgadas
\end{itemize}

\textbf{Interpretación}: igual que Pearson ($-1$ a $1$), pero mide asociación monotónica.

\vspace{0.3em}
\begin{exampleblock}{Ejemplos}
Estrato (ordinal: 1, ..., 6) vs puntaje de ingl\'es: Spearman es m\'as apropiado que Pearson porque estrato es ordinal.
Rating de un restaurante en Rappi (1--5 estrellas, ordinal) vs.\ precio del plato $\to$ Spearman es m\'as apropiado que Pearson.
\end{exampleblock}
\end{frame}

\begin{frame}[fragile]{Correlación de Spearman en R}
\begin{lstlisting}
# Correlacion de Spearman
cor(icfes$FAMI_ESTRATOVIVIENDA, icfes$MOD_INGLES_PUNT,
    method = "spearman", use = "complete.obs")
# [1] 0.428

# Comparar Pearson vs Spearman
cat("Pearson:", cor(icfes$FAMI_ESTRATOVIVIENDA,
                    icfes$MOD_INGLES_PUNT,
                    use = "complete.obs"), "\n")
cat("Spearman:", cor(icfes$FAMI_ESTRATOVIVIENDA,
                     icfes$MOD_INGLES_PUNT,
                     method = "spearman",
                     use = "complete.obs"), "\n")
\end{lstlisting}
\end{frame}

\begin{frame}[fragile]{Matriz de correlación de Spearman}
\begin{lstlisting}
# Matriz de correlacion de Spearman
matriz_cor_spearman <- cor(variables_interes,
                            method = "spearman",
                            use = "complete.obs")

corrplot(matriz_cor_spearman, method = "circle",
         type = "upper", addCoef.col = "black")
\end{lstlisting}

Comparar las matrices de Pearson y Spearman permite detectar relaciones no lineales.
\end{frame}

\begin{frame}[fragile]{Ejemplo: correlación entre módulos Saber Pro}
\begin{lstlisting}
# Seleccionar los 5 modulos genericos (con alias cortos)
modulos <- icfes %>%
  select(Ingles = MOD_INGLES_PUNT,
         Lectura = MOD_LECTURA_CRITICA_PUNT,
         Cuant = MOD_RAZONA_CUANTITAT_PUNT,
         Ciudad = MOD_COMPETEN_CIUDADA_PUNT,
         Escrit = MOD_COMUNICACION_ESCRITA_PUNT)

# Matriz de correlacion
round(cor(modulos, use = "complete.obs"), 3)

# Visualizacion avanzada con GGally
library(GGally)
ggpairs(modulos,
        title = "Correlacion - Modulos Saber Pro")
\end{lstlisting}
\end{frame}

% ============================================================================
\section{Paradoja de Simpson}
% ============================================================================

\begin{frame}{Paradoja de Simpson}
\begin{block}{Definición}
La \textbf{Paradoja de Simpson} ocurre cuando una tendencia que aparece en varios grupos \alert{se invierte} cuando los grupos se combinan.
\end{block}

\textbf{Implicaciones}:
\begin{itemize}
\item Una relación puede cambiar de dirección al agregar o desagregar datos
\item Controlando por una tercera variable (confusora), la relación puede invertirse
\item Muestra la importancia de identificar variables confusoras
\item Fundamental para el diseño de experimentos y análisis causal
\end{itemize}

\vspace{0.5em}
\begin{exampleblock}{Ejemplo cotidiano}
Una app de delivery puede tener mejor rating \textbf{agregado} que otra, pero peor rating en cada categor\'ia individual (comida, rapidez, empaque) --- porque una app recibe m\'as rese\~nas en categor\'ias dif\'iciles.
\end{exampleblock}

\begin{alertblock}{Lecci\'on clave}
\alert{Siempre explorar subgrupos antes de generalizar}. Lo que es verdad para el agregado puede no ser verdad para los componentes (y viceversa).
\end{alertblock}
\end{frame}

\begin{frame}{Ejemplo clásico: admisiones Berkeley 1973}
\textbf{Contexto}: demanda por discriminación de género en admisiones de posgrado.

\vspace{0.5em}
\textbf{Datos agregados}:
\begin{table}
\centering
\begin{tabular}{lcc}
\toprule
& Admitidos & Tasa de admisión \\
\midrule
Hombres & 1198 / 2691 & 44\% \\
Mujeres & 557 / 1835 & 30\% \\
\bottomrule
\end{tabular}
\end{table}

Conclusión agregada: parece haber sesgo contra mujeres (44\% vs 30\%).

\vspace{0.5em}
\textbf{Datos desagregados por departamento}:
\begin{itemize}
\item En casi todos los departamentos, las mujeres tenían tasas de admisión \alert{iguales o mayores}
\item Pero las mujeres aplicaban más a departamentos muy competitivos (baja tasa de admisión)
\item Los hombres aplicaban más a departamentos menos competitivos (alta tasa de admisión)
\end{itemize}

\alert{Variable confusora}: competitividad del departamento
\end{frame}

\begin{frame}[fragile]{Ejemplo con datos ICFES (1/2)}
\begin{lstlisting}
# Ejemplo simulado: relacion entre estrato y puntaje
# agregada vs desagregada por tipo de IES

# Correlacion agregada (todos los estudiantes juntos)
cor_agregada <- cor(icfes$FAMI_ESTRATOVIVIENDA,
                    icfes$MOD_INGLES_PUNT,
                    use = "complete.obs")

# Correlacion por tipo de IES
cor_desagregada <- icfes %>%
  group_by(INST_CARACTER_ACADEMICO) %>%
  summarise(
    n = n(),
    cor = cor(FAMI_ESTRATOVIVIENDA, MOD_INGLES_PUNT,
              use = "complete.obs")
  )
\end{lstlisting}
\end{frame}

\begin{frame}[fragile]{Ejemplo con datos ICFES (2/2)}
\begin{lstlisting}
# Scatter plot mostrando la paradoja
ggplot(icfes, aes(x = FAMI_ESTRATOVIVIENDA, y = MOD_INGLES_PUNT,
                  color = INST_CARACTER_ACADEMICO)) +
  geom_point(alpha = 0.3) +
  geom_smooth(method = "lm", se = FALSE) +  # Lineas por grupo
  geom_smooth(aes(group = 1), method = "lm", color = "black",
              linetype = "dashed", se = FALSE) +  # Linea agregada
  labs(title = "Paradoja de Simpson: estrato vs puntaje") +
  theme_minimal()
\end{lstlisting}
\end{frame}

\begin{frame}{Por qué la regresión múltiple es necesaria}
\begin{block}{Problema fundamental}
Las correlaciones bivariadas pueden ser engañosas:
\begin{itemize}
\item No controlan por variables confusoras
\item No separan efectos de múltiples factores simultáneos
\item Pueden reflejar asociaciones espurias
\end{itemize}
\end{block}

\textbf{Solución: Regresión múltiple}
\begin{itemize}
\item Estima el efecto de una variable \alert{manteniendo otras constantes}
\item Controla por variables confusoras
\item Ejemplo: efecto de Saber 11 en Saber Pro, \alert{controlando por} tipo de IES, estrato, género
\end{itemize}

\vspace{0.3em}
\begin{alertblock}{Adelanto}
En sesiones 5--6 estudiaremos regresión lineal simple y múltiple.
\end{alertblock}
\end{frame}

% ============================================================================
\section{Pipeline completo de EDA}
% ============================================================================

\begin{frame}{Flujo de trabajo EDA}
\begin{block}{Análisis Exploratorio de Datos (EDA)}
Proceso iterativo para entender la estructura, patrones y anomal\'ias en los datos \alert{antes} de modelar. Es como cuando stalkean un perfil de Instagram: primero ven el feed general, luego las fotos, luego los comments, luego los tagged\ldots\ El EDA hace lo mismo con datos.
\end{block}

\textbf{Pipeline t\'ipico}:
\begin{enumerate}
\item \textbf{Cargar}: importar datos, verificar dimensiones
\item \textbf{Limpiar}: identificar y tratar valores faltantes, duplicados, errores
\item \textbf{Transformar}: crear variables derivadas, recodificar, filtrar
\item \textbf{Resumir}: estadísticas descriptivas (media, mediana, sd, etc.)
\item \textbf{Visualizar}: histogramas, box plots, scatter plots, correlaciones
\item \textbf{Segmentar}: análisis por subgrupos (tipo IES, estrato, región)
\item \textbf{Identificar anomalías}: outliers, valores extremos, inconsistencias
\item \textbf{Documentar}: hallazgos clave, preguntas para investigar
\end{enumerate}
\end{frame}

\begin{frame}[fragile]{EDA completo de ICFES NI: cargar y limpiar}
\small
\begin{lstlisting}
library(tidyverse)
library(moments)
library(corrplot)

# 1. CARGAR
icfes <- read_csv("datos/saber_pro_ni_2024.csv")
dim(icfes)  # dimension
glimpse(icfes)  # estructura

# 2. LIMPIAR
# Valores faltantes
colSums(is.na(icfes))
# Eliminar filas con NA en variable clave
icfes_clean <- icfes %>%
  filter(!is.na(MOD_INGLES_PUNT))

# Duplicados
icfes_clean <- icfes_clean %>% distinct()
\end{lstlisting}
\end{frame}

\begin{frame}[fragile]{EDA completo de ICFES NI: transformar}
\begin{lstlisting}
# 3. TRANSFORMAR
icfes_clean <- icfes_clean %>%
  mutate(
    # Categoria de desempeno en ingles
    nivel_ingles = cut(MOD_INGLES_PUNT,
                       breaks = c(0, 140, 160, 180, 200),
                       labels = c("Bajo", "Medio", "Alto", "Superior")),
    # Promedio de modulos genericos
    promedio_modulos = (MOD_INGLES_PUNT + MOD_LECTURA_CRITICA_PUNT +
                        MOD_RAZONA_CUANTITAT_PUNT) / 3
  )
\end{lstlisting}
\end{frame}

\begin{frame}[fragile]{EDA completo: resumir}
\begin{lstlisting}
# 4. RESUMIR
summary(icfes_clean$MOD_INGLES_PUNT)
sd(icfes_clean$MOD_INGLES_PUNT)
skewness(icfes_clean$MOD_INGLES_PUNT)
\end{lstlisting}
\end{frame}

\begin{frame}[fragile]{EDA completo: visualizar (histograma y box plot)}
\begin{lstlisting}
# 5. VISUALIZAR
# Histograma
ggplot(icfes_clean, aes(x = MOD_INGLES_PUNT)) +
  geom_histogram(bins = 30, fill = "steelblue") +
  theme_minimal()

# Box plot por tipo IES
ggplot(icfes_clean, aes(x = INST_CARACTER_ACADEMICO,
                        y = MOD_INGLES_PUNT)) +
  geom_boxplot() + theme_minimal()
\end{lstlisting}
\end{frame}

\begin{frame}[fragile]{EDA completo: visualizar (scatter y correlaciones)}
\begin{lstlisting}
# Scatter plot: Saber 11 vs Saber Pro
ggplot(icfes_clean, aes(x = PUNT_INGLES_S11, y = MOD_INGLES_PUNT)) +
  geom_point(alpha = 0.3) +
  geom_smooth(method = "lm") +
  theme_minimal()

# Matriz de correlacion
cor_matrix <- cor(icfes_clean %>%
  select(starts_with("MOD_")), use = "complete.obs")
corrplot(cor_matrix, method = "color")
\end{lstlisting}
\end{frame}

\begin{frame}[fragile]{EDA completo: segmentar}
\begin{lstlisting}
# 6. SEGMENTAR
# Por tipo de IES
icfes_clean %>%
  group_by(INST_CARACTER_ACADEMICO) %>%
  summarise(
    n = n(),
    media = mean(MOD_INGLES_PUNT),
    mediana = median(MOD_INGLES_PUNT),
    sd = sd(MOD_INGLES_PUNT),
    asimetria = skewness(MOD_INGLES_PUNT)
  )

# Por estrato
icfes_clean %>%
  filter(!is.na(FAMI_ESTRATOVIVIENDA)) %>%
  group_by(FAMI_ESTRATOVIVIENDA) %>%
  summarise(across(MOD_INGLES_PUNT,
                   list(media = mean, mediana = median, sd = sd)))
\end{lstlisting}
\end{frame}

\begin{frame}[fragile]{EDA completo: identificar anomalías}
\begin{lstlisting}
# 7. IDENTIFICAR ANOMALIAS
# Outliers (metodo IQR)
Q1 <- quantile(icfes_clean$MOD_INGLES_PUNT, 0.25)
Q3 <- quantile(icfes_clean$MOD_INGLES_PUNT, 0.75)
IQR_val <- Q3 - Q1
outliers <- icfes_clean %>%
  filter(MOD_INGLES_PUNT < Q1 - 1.5*IQR_val |
         MOD_INGLES_PUNT > Q3 + 1.5*IQR_val)
nrow(outliers)
\end{lstlisting}
\end{frame}

\begin{frame}{Checklist EDA}
\begin{alertblock}{Lista de verificación antes de modelar}
\begin{itemize}
\item[$\square$] ¿Entiendo qué mide cada variable?
\item[$\square$] ¿Cuántos valores faltantes hay? ¿Por qué faltan?
\item[$\square$] ¿Hay duplicados?
\item[$\square$] ¿Hay outliers? ¿Son errores o datos reales?
\item[$\square$] ¿Cuál es la distribución de cada variable (forma, centro, dispersión)?
\item[$\square$] ¿Hay relaciones entre variables? ¿Cuáles son las más fuertes?
\item[$\square$] ¿Hay diferencias importantes entre subgrupos?
\item[$\square$] ¿Hay variables confusoras que deba controlar?
\item[$\square$] ¿Los datos son representativos de la población de interés?
\item[$\square$] ¿Qué preguntas sustantivas puedo responder con estos datos?
\end{itemize}
\end{alertblock}

\vspace{0.5em}
\alert{Documentar hallazgos en un reporte o Rmarkdown!}
\end{frame}

% ============================================================================
\section{Resumen}
% ============================================================================

\begin{frame}{Resumen de la sesión}
\textbf{Conceptos clave cubiertos}:

\begin{enumerate}
\item \textbf{Distribuciones de frecuencia}: absolutas, relativas, acumuladas
\item \textbf{Visualización univariada}: histograma, density plot, box plot, barras
\item \textbf{Visualización bivariada}: scatter plot, heatmap, tablas cruzadas
\item \textbf{Covarianza y correlación}: Pearson (lineal), Spearman (monotónica)
\item \textbf{Paradoja de Simpson}: por qué desagregar es crucial
\item \textbf{Pipeline EDA}: cargar → limpiar → explorar → visualizar → reportar
\end{enumerate}

\vspace{1em}
\begin{alertblock}{Mensajes clave}
\begin{itemize}
\item \alert{Visualizar siempre antes de resumir numéricamente}
\item \alert{Correlación $\neq$ causalidad}
\item \alert{Explorar subgrupos antes de agregar}
\item \alert{El EDA es iterativo, no lineal}
\end{itemize}
\end{alertblock}
\end{frame}

\begin{frame}{Próxima sesión: Probabilidad, Normal y Muestreo}
\textbf{Vie 21 ago 2026}

\textbf{Temas a cubrir}:
\begin{itemize}
\item Variables aleatorias y distribuciones de probabilidad
\item Distribución normal: propiedades, cálculo de probabilidades, cuantiles
\item Distribuciones t, chi-cuadrado y F
\item Muestreo y distribución muestral de $\bar{X}$
\item \alert{Teorema Central del Límite} (fundamento de la inferencia)
\item Propiedades de estimadores
\end{itemize}

\vspace{1em}
\textbf{Preparación}:
\begin{itemize}
\item Repasar conceptos de probabilidad básica (si es necesario)
\item Instalar paquete \paquete{moments} en R
\item Familiarizarse con funciones \code{dnorm()}, \code{pnorm()}, \code{qnorm()}
\end{itemize}
\end{frame}

\end{document}
